% 建模活动单
% 三个活动分别用指数、对数、幂三种模型，走完整的建模过程。
\section{活动一：人口增长——指数模型}

\subsection{情境}
某城市 2000 年人口 800 万，2010 年 1000 万，2020 年 1250 万。

\subsection{提出问题}
按这个趋势，2050 年大约多少人口？哪一年超过 2000 万？

\subsection{选模型}
人口增长常用指数模型
\[
  P(t)=P_0a^t\quad(a>1),
\]
其中 $t$ 是从 2000 年起算的年数。
为什么不用二次函数？因为人口增长是“按当前规模成比例增长”，
即增长量与现有规模成正比，而不是与时间的平方成正比。

\subsection{拟合}
由 $P(0)=800$、$P(10)=1000$ 得 $a^{10}=\dfrac{1000}{800}=1.25$，
故 $a=1.25^{1/10}\approx1.0226$。
检验：$P(20)=800\times1.0226^{20}\approx1250$（万），与实际吻合。

\subsection{预测}
\[
  P(50)=800\times1.0226^{50}\approx800\times3.05\approx2440\text{（万）}.
\]
超过 $2000$ 万：解 $800\times1.0226^t=2000$，即 $1.0226^t=2.5$，
\[
  t=\frac{\ln2.5}{\ln1.0226}\approx41.3,
\]
即大约在 2041 年超过 2000 万。

\subsection{检验与反思}
指数模型在短期内吻合，但长期看它会无限增长，这不合理——资源是有限的。
更真实的模型是逻辑斯蒂模型（选学）：增长到一定规模后会趋于饱和。
\key{建模的第一步是选对模型，选错了，计算再准也没用。}

\section{活动二：声音强度——对数模型}

\subsection{情境}
声音的强度 $I$ 相差极大：耳语约 $10^{-10}\ \mathrm{W/m^2}$，摇滚乐约 $10^{-1}\ \mathrm{W/m^2}$。

\subsection{提出问题}
怎样用一个“好写好看”的数来描述声音的大小？

\subsection{选模型}
用分贝
\[
  L=10\lg\frac{I}{I_0},\qquad I_0=10^{-12}\ \mathrm{W/m^2}.
\]
为什么用对数？因为人耳对强度的感受，近似与强度的对数成正比。

\subsection{计算}
耳语：$L=10\lg\dfrac{10^{-10}}{10^{-12}}=10\lg100=20$（dB）。
摇滚乐：$L=10\lg\dfrac{10^{-1}}{10^{-12}}=10\lg10^{11}=110$（dB）。

\subsection{检验与反思}
强度相差 $10^9$ 倍，分贝只相差 $90$——对数把“倍数”变成了“差值”。
同样的思路：里氏震级 $M=\lg\dfrac{A}{A_0}$、溶液 $\mathrm{pH}=-\lg c(\mathrm{H}^+)$。
\key{凡是要描述“跨好几个数量级”的量，就考虑取对数。}

\section{活动三：行星轨道——幂模型}

\subsection{情境}
行星绕太阳公转的周期 $T$ 与轨道半长轴 $r$ 之间有什么关系？

\subsection{数据}
\begin{center}
\footnotesize
\setlength{\tabcolsep}{4pt}
\renewcommand{\arraystretch}{1.4}
\begin{tabular}{p{2.4cm}p{3.6cm}p{3.6cm}}
\toprule
行星 & 半长轴 $r$（AU） & 周期 $T$（年） \\
\midrule
地球 & $1$ & $1$ \\
火星 & $1.52$ & $1.88$ \\
木星 & $5.20$ & $11.86$ \\
\bottomrule
\end{tabular}
\end{center}

\subsection{提出问题}
已知 $r$，怎样求 $T$？

\subsection{选模型与拟合}
试幂模型 $T=kr^{\alpha}$。
由地球数据得 $k=1$。由火星数据：
\[
  1.88=1.52^{\alpha}\implies\alpha=\frac{\ln1.88}{\ln1.52}\approx1.497\approx\frac32.
\]
检验木星：$T=5.20^{3/2}\approx11.86$（年），与实际吻合。

\subsection{结论}
\[
  T=r^{3/2}\quad\Longleftrightarrow\quad T^2=r^3,
\]
这正是开普勒第三定律。幂指数 $\alpha=\dfrac32$ 介于指数增长与对数增长之间。

\subsection{检验与反思}
三种模型各有适用场景：
\begin{itemize}
  \item 指数模型 $y=a^x$：按当前规模成比例增长（人口、复利、细菌）；
  \item 对数模型 $y=\log_a x$：跨数量级的量（分贝、震级、pH）；
  \item 幂模型 $y=x^{\alpha}$：由物理定律给出的幂律关系（开普勒定律）。
\end{itemize}

\section{建模报告要求}
每个活动写一份 300 字左右的报告，包含四部分：
\begin{enumerate}
  \item 要解决的问题（用一两句话说清楚）；
  \item 选的模型与理由（为什么选它，而不是别的）；
  \item 拟合过程与检验（代了哪些数据，误差多大）；
  \item 预测与反思（模型在什么情况下会失效）。
\end{enumerate}
\key{评价标准}：模型选得对不对，比算得准不准更重要。
